\section{Agentic Security Risk Map (OWASP Top 10 for Agents)}

The rapid integration of AI agents into production environments introduces a novel attack surface, necessitating a dedicated security framework. Just as the OWASP Top 10 has guided web application security for years, a similar mapping is crucial for agentic applications \cite{OWASP2026}. These risks stem from the unique characteristics of AI agents, such as their reliance on large language models (LLMs), their ability to interact with external tools, and their dynamic, often unpredictable, behavior \cite{Reger2026}. Understanding and mitigating these risks is paramount to deploying AI agents safely and effectively.

Here, we outline key security risks for production AI agents, drawing parallels and distinctions with established security principles:

1.  \textbf{Goal Hijacking (Prompt Injection \& Manipulation):}     \textit{   \textbf{Description:} This is arguably the most prevalent and unique threat to AI agents. It occurs when an attacker manipulates an agent's instructions or goals through carefully crafted inputs, causing the agent to perform unintended actions or reveal sensitive information \cite{Reger2026}. This can range from subtle modifications of user prompts to direct injection of malicious commands within data the agent processes.     }   \textbf{Mechanism:} Attackers exploit the agent's reliance on natural language processing and its instruction-following capabilities. For example, an agent designed to summarize documents could be tricked into ignoring its original instruction and instead executing malicious code embedded within the document it is supposed to summarize \cite{OWASP2026}.     *   \textbf{Mitigation:} Robust input validation and sanitization are critical. Employing techniques like prompt shielding, instruction defense, and using separate LLMs for processing user input versus internal instructions can help \cite{Reger2026}. Sandboxing the agent's execution environment is also vital (see Risk 3).

2.  \textbf{Tool Misuse (Unauthorized Access \& Action):}     \textit{   \textbf{Description:} AI agents often leverage external tools (APIs, databases, scripts) to interact with the environment. This risk arises when an agent is tricked or compelled into misusing these tools, leading to data breaches, unauthorized system modifications, or denial of service \cite{Reger2026}.     }   \textbf{Mechanism:} Similar to Goal Hijacking, attackers can manipulate the agent's intent to call tools with malicious parameters or to access sensitive tools they shouldn't have access to. An agent tasked with sending emails might be manipulated to send phishing messages or spam \cite{OWASP2026}.     *   \textbf{Mitigation:} Implement a strict least-privilege access control model for all tools accessible by agents. Tools should have clearly defined capabilities and inputs/outputs, with rigorous validation before execution. An agent's access to tools should be dynamically assessed based on its current, verified goal \cite{Reger2026}.

3.  \textbf{Insecure Output Handling (Data Leakage \& Malicious Content):}     \textit{   \textbf{Description:} Agents process and generate outputs that can be sensitive or, if malicious, pose a risk to downstream systems or users. This risk covers scenarios where agents leak confidential data in their responses or generate harmful content (e.g., biased text, malicious code snippets) \cite{OWASP2026}.     }   \textbf{Mechanism:} An agent might inadvertently include sensitive information from its training data or from retrieved documents in its output. Conversely, a compromised agent or one susceptible to prompt injection could be made to generate harmful code or misinformation \cite{Reger2026}.     *   \textbf{Mitigation:} Implement output filtering and sanitization. Use content moderation models to scan agent outputs for sensitive information, harmful content, or malicious code before they are presented to users or processed by other systems. Carefully manage the training data to minimize inherent biases and sensitive data leakage \cite{Reger2026}.

4.  \textbf{Training Data Poisoning:}     \textit{   \textbf{Description:} The performance and behavior of an AI agent are heavily dependent on its training data. This risk involves an attacker subtly manipulating the training dataset to introduce vulnerabilities, biases, or backdoors into the agent's model \cite{Reger2026}.     }   \textbf{Mechanism:} Attackers might inject malicious examples into datasets used for fine-tuning LLMs or embedding knowledge bases. This could cause the agent to consistently provide incorrect information, exhibit discriminatory behavior, or even execute malicious code when specific conditions are met \cite{OWASP2026}.     *   \textbf{Mitigation:} Rigorous data validation, integrity checks, and provenance tracking for all training datasets are essential. Employ outlier detection and anomaly detection techniques during the training process. Regularly audit and re-validate models against known good datasets \cite{Reger2026}.

5.  \textbf{Insecure Sandboxing \& Environment Escape:}     \textit{   \textbf{Description:} Agents, especially code-executing agents, often run within sandboxed environments for security. This risk occurs when an agent or an attacker exploiting the agent can break out of the sandbox, gaining unauthorized access to the host system or other isolated environments \cite{Reger2026}.     }   \textbf{Mechanism:} Vulnerabilities in the sandboxing technology itself, misconfigurations, or sophisticated exploits within the agent's code can allow it to escape its intended confines \cite{OWASP2026}.     *   \textbf{Mitigation:} Utilize robust, up-to-date sandboxing technologies (e.g., containers, virtual machines) with minimal privilege configurations. Continuously monitor sandbox activity for suspicious behavior and implement network segmentation to limit the blast radius of any potential escape \cite{Reger2026}.

6.  \textbf{Over-reliance \& Unverified Automation:}     \textit{   \textbf{Description:} A critical human factor risk where users or systems place undue trust in an agent's output or actions without proper verification, especially for high-stakes decisions \cite{Reger2026}.     }   \textbf{Mechanism:} The perceived intelligence and efficiency of AI agents can lead to a passive acceptance of their results, bypassing necessary human oversight or validation steps. This can result in critical errors going unnoticed \cite{OWASP2026}.     *   \textbf{Mitigation:} Design agents with explicit human-in-the-loop checkpoints for critical tasks. Provide clear indicators of confidence levels for agent outputs. Educate users on the limitations of AI and the importance of verification \cite{Reger2026}.

7.  \textbf{Model Inversion \& Membership Inference:}     \textit{   \textbf{Description:} These attacks aim to infer sensitive information about the agent's training data. Model inversion attempts to reconstruct training data samples, while membership inference tries to determine if a specific data record was part of the training set \cite{Reger2026}.     }   \textbf{Mechanism:} By querying the agent with specific inputs and analyzing its outputs, attackers can potentially deduce patterns or even specific data points used during training, which might include personally identifiable information (PII) or proprietary business data \cite{OWASP2026}.     *   \textbf{Mitigation:} Employ privacy-preserving techniques during training, such as differential privacy. Limit the granularity and specificity of information returned by the agent, especially when dealing with sensitive datasets \cite{Reger2026}.

8.  \textbf{Denial of Service (DoS) \& Resource Exhaustion:}     \textit{   \textbf{Description:} Agents, particularly those involving complex LLM computations or extensive tool usage, can be resource-intensive. Attackers can trigger agents in ways that intentionally exhaust computational resources, leading to service degradation or unavailability \cite{Reger2026}.     }   \textbf{Mechanism:} Tricking an agent into performing extremely computationally expensive tasks, making recursive calls, or overwhelming its tool integrations can lead to DoS conditions \cite{OWASP2026}.     *   \textbf{Mitigation:} Implement rate limiting, resource quotas, and circuit breakers for agent operations and tool usage. Optimize agent execution paths and LLM calls for efficiency. Monitor resource utilization closely \cite{Reger2026}.

9.  \textbf{Supply Chain Vulnerabilities:}     \textit{   \textbf{Description:} AI agents often rely on numerous third-party libraries, pre-trained models, and external services. Vulnerabilities in any of these components can compromise the entire agent system \cite{Reger2026}.     }   \textbf{Mechanism:} An attacker might compromise a popular LLM library, a pre-trained model on a public repository, or an external API service that the agent interacts with, introducing malicious code or data \cite{OWASP2026}.     *   \textbf{Mitigation:} Maintain a strict software bill of materials (SBOM) for all agent components. Vet all third-party dependencies, use trusted sources for models and libraries, and implement security scanning for dependencies \cite{Reger2026}.

10. \textbf{Lack of Agent Observability \& Traceability:}     \textit{   \textbf{Description:} Complex agentic systems can be difficult to monitor and audit. A lack of clear visibility into agent decision-making, actions, and interactions makes it hard to detect, diagnose, and respond to security incidents \cite{Reger2026}.     }   \textbf{Mechanism:} Without adequate logging and tracing, it's challenging to understand how an agent arrived at a particular decision, which tools it used, or what data it accessed, hindering incident response and forensic analysis \cite{OWASP2026}.     *   \textbf{Mitigation:} Implement comprehensive logging and tracing for all agent activities, including prompt processing, tool calls, LLM interactions, and outputs. Establish clear audit trails for security-sensitive operations \cite{Reger2026}.

Addressing these agent-specific security risks requires a proactive and layered security strategy, integrating principles of secure coding, robust access control, continuous monitoring, and comprehensive validation throughout the agent lifecycle \cite{Reger2026}.
